\SourcePage{547}{550}
\section{The Schrödinger equation in a magnetic field}

A particle with spin also possesses an ``intrinsic'' magnetic moment
$\boldsymbol\mu$. Its quantum-mechanical operator is proportional to the spin
operator $\hat{\mathbf s}$ and may therefore be written
\begin{equation}
 \hat{\boldsymbol\mu}=\frac{\mu}{s}\hat{\mathbf s},
 \tag{111.1}
\end{equation}
where $s$ is the magnitude of the particle's spin and $\mu$ is a constant
characteristic of the particle. The eigenvalues of the projection of the
magnetic moment are $\mu_z=\mu\sigma/s$. Thus the coefficient $\mu$, which is
usually referred to simply as the magnitude of the magnetic moment, is the
largest possible value of $\mu_z$ and is attained for the spin projection
$\sigma=s$.

The ratio $\mu/(\hbar s)$ gives the ratio of the particle's intrinsic magnetic
moment to its intrinsic mechanical angular momentum when both are directed
along the $z$ axis. For ordinary orbital angular momentum this ratio is, as is
known, $e/(2mc)$ (see II, \S~44). The proportionality coefficient between the
intrinsic magnetic moment and the spin of a particle is different. For the
electron it is $-|e|/mc$, twice the ordinary value (this value follows
theoretically from the relativistic Dirac wave equation; see IV, \S~33). The
intrinsic magnetic moment of the electron, whose spin is $1/2$, is consequently
$-\mu_B$, where
\begin{equation}
 \mu_B=\frac{|e|\hbar}{2mc}=0.927\cdot10^{-20}\frac{\text{erg}}{\text{G}}.
 \tag{111.2}
\end{equation}
This quantity is called the \emph{Bohr magneton}.

The magnetic moments of heavy particles are conventionally measured in
\emph{nuclear magnetons}, defined as $e\hbar/(2m_pc)$, where $m_p$ is the
proton mass. Experiment gives the intrinsic magnetic moment of the proton as
$2.79$ nuclear magnetons, directed along the spin. The neutron's magnetic
moment is directed opposite to its spin and has magnitude $1.91$ nuclear
magnetons.

\SourcePage{548}{551}
Notice that the quantities $\boldsymbol\mu$ and $\mathbf s$ on the two sides
of (111.1) have, as they must, the same vector character: both are axial
vectors. An analogous equality for the electric dipole moment $\mathbf d$,
namely $\mathbf d=\mathrm{const}\cdot\mathbf s$, would be incompatible with
symmetry under inversion of the coordinates, since inversion would change the
relative sign of the two sides\footnote{It should be noted that this relation (and hence the existence of an electric moment in an elementary particle) would also be incompatible with time-reversal symmetry: reversing the sign of time leaves $\mathbf d$ unchanged but reverses the spin (as is clear, for example, from the definitions of these quantities for orbital motion: only coordinates enter the definition of $\mathbf d$, whereas the definition of the moment also contains the particle's velocity).}
.

In nonrelativistic quantum mechanics a magnetic field can be treated only as
an external field. The magnetic interaction between particles is a
relativistic effect, and a consistent relativistic theory is required to take
it into account.

In classical theory the Hamilton function of a charged particle in an
electromagnetic field is
\[
 H=\frac{1}{2m}\left(\mathbf p-\frac ec\mathbf A\right)^2+e\varphi,
\]
where $\varphi$ and $\mathbf A$ are respectively the scalar and vector
potentials of the field, while $\mathbf p$ is the generalized momentum of the
particle (see II, \S~16). If the particle has no spin, the transition to
quantum mechanics is made in the usual way: the generalized momentum is
replaced by the operator $\hat{\mathbf p}=-i\hbar\boldsymbol\nabla$, giving the
Hamiltonian\footnote{Here we use the same letter $\mathbf p$ for the
generalized momentum as for the ordinary momentum, instead of the $\mathbf P$
used in II, \S~16, in order to emphasize that the same operator corresponds to
it.}
\begin{equation}
 \hat H=\frac{1}{2m}\left(\hat{\mathbf p}-\frac ec\mathbf A\right)^2+e\varphi.
 \tag{111.3}
\end{equation}

If the particle has spin, however, this procedure is insufficient. Its
intrinsic magnetic moment interacts directly with the magnetic field. This
interaction is entirely absent from the classical Hamilton function, since
spin itself, being a purely quantum effect, disappears in the classical
limit. The correct Hamiltonian is obtained by adding to (111.3) the term
$-\hat{\boldsymbol\mu}\mathbf H$, corresponding to the energy of a magnetic
moment $\boldsymbol\mu$ in a field $\mathbf H$. Thus the
\SourcePage{549}{552}
Hamiltonian of a particle with spin is\footnote{Using the same letter for the
magnetic field and the Hamiltonian cannot cause confusion: the symbol for the
Hamiltonian carries a hat.}
\begin{equation}
 \hat H=\frac{1}{2m}\left(\hat{\mathbf p}-\frac ec\mathbf A\right)^2-\hat{\boldsymbol\mu}\mathbf H+e\varphi.
 \tag{111.4}
\end{equation}
When expanding the square $(\hat{\mathbf p}-(e/c)\mathbf A)^2$, one must bear
in mind that the operator $\hat{\mathbf p}$ does not in general commute with
the vector $\mathbf A$, which is a function of the coordinates. We must
therefore write
\begin{equation}
 \hat H=\frac{1}{2m}\hat{\mathbf p}^{2}-\frac{e}{2mc}(\hat{\mathbf p}\mathbf A+\mathbf A\hat{\mathbf p})+\frac{e^2}{2mc^2}\mathbf A^2-\frac{\mu}{s}\hat{\mathbf s}\mathbf H+e\varphi.
 \tag{111.5}
\end{equation}
The commutation rule (16.4) for the momentum operator and an arbitrary
function of the coordinates gives
\begin{equation}
 \hat{\mathbf p}\mathbf A-\mathbf A\hat{\mathbf p}=-i\hbar\Div{\mathbf A}.
 \tag{111.6}
\end{equation}

Thus $\hat{\mathbf p}$ and $\mathbf A$ commute if $\Div{\mathbf A}=0$. This
condition holds, in particular, for a uniform field when its vector potential
is chosen as
\begin{equation}
 \mathbf A=(1/2)(\mathbf H\times\mathbf r).
 \tag{111.7}
\end{equation}

The equation $i\hbar\,\partial\Psi/\partial t=\hat H\Psi$, with Hamiltonian
(111.4), is the generalization of the Schrödinger equation to the presence of
a magnetic field. The wave functions acted on by the Hamiltonian in this
equation are symmetric spinors of rank $2s$.

Wave functions of a particle in an electromagnetic field have an ambiguity
associated with the ambiguity in the field potentials. As is known (see II,
\S~18), the latter are determined only up to a \emph{gauge transformation}
\begin{equation}
 \mathbf A\to\mathbf A+\boldsymbol\nabla f,\qquad
 \varphi\to\varphi-\frac1c\frac{\partial f}{\partial t},
 \tag{111.8}
\end{equation}
where $f$ is an arbitrary function of the coordinates and time. Such a
transformation does not affect the field strengths. It must therefore leave
the solutions of the wave equation essentially unchanged as well; in
particular, $|\Psi|^2$ must remain invariant. It is indeed readily verified
that the original equation is recovered if, together with the replacement
(111.8) in the Hamiltonian, the wave function is replaced according to
\begin{equation}
 \Psi\to\Psi\exp\left(\frac{ie}{\hbar c}f\right).
 \tag{111.9}
\end{equation}

\SourcePage{550}{553}
This ambiguity of the wave function has no effect on any physically
meaningful quantity whose definition does not explicitly contain the
potentials.

In classical mechanics a particle's generalized momentum is related to its velocity by $m\mathbf v=\mathbf p-e\mathbf A/c$. To find the operator $\hat{\mathbf v}$ in quantum mechanics, one must form the commutator of the vector $\mathbf r$ with the Hamiltonian. A straightforward calculation gives

\begin{equation}
 m\hat{\mathbf v}=\hat{\mathbf p}-\frac ec\mathbf A,
 \tag{111.10}
\end{equation}
which is exactly analogous to the classical relation. The operators of the
velocity components obey the commutation rules
\begin{equation}
 \left\{\hat v_x,\hat v_y\right\}=i\frac{e\hbar}{m^2c}H_z,\qquad
 \left\{\hat v_y,\hat v_z\right\}=i\frac{e\hbar}{m^2c}H_x,\qquad
 \left\{\hat v_z,\hat v_x\right\}=i\frac{e\hbar}{m^2c}H_y,
 \tag{111.11}
\end{equation}
These relations are readily verified by direct calculation. Thus, in a magnetic field the operators for the three velocity components of a (charged) particle do not commute. Consequently, the particle cannot simultaneously have definite velocity values along all three directions.


For motion in a magnetic field, time-reversal symmetry holds only if the sign
of the field $\mathbf H$, and of the vector potential $\mathbf A$, is reversed
as well. This means (see \S\S~18 and 60) that the Schrödinger equation
$\hat H\psi=E\psi$ must retain its form upon passage to the complex-conjugate
quantities accompanied by reversal of the sign of $\mathbf H$. This is
immediately evident for every term in the Hamiltonian (111.4) except
$-\hat{\mathbf s}\mathbf H$. Under the specified transformation the term
$-\hat{\mathbf s}\mathbf H\psi$ in the Schrödinger equation becomes
$\hat{\mathbf s}^{*}\mathbf H\psi^{*}$ and appears at first sight to violate
the required invariance, since the operator $\hat{\mathbf s}^{*}$ is not equal
to $-\hat{\mathbf s}$.

One must take into account, however, that the wave function is actually a
contravariant spinor $\psi^{\lambda\mu\ldots}$, which on complex conjugation
becomes the covariant spinor $\psi^{\lambda\mu\ldots *}$ (see \S~60). The
spinor $\psi^{*}_{\lambda\mu\ldots}$, on the other hand, is contravariant.
Using definitions (57.4) and (57.5) to find the components of
$(\hat{\mathbf s}\mathbf H\psi)^*$ and expressing them in terms of
$\psi^{*}_{\lambda\mu\ldots}$, one verifies that time reversal leads to a
Schrödinger equation for the components $\psi^{*}_{\lambda\mu\ldots}$ of the
same form as the original equation for the components
$\psi^{\lambda\mu\ldots}$.
